\documentclass[11pt]{article}
\usepackage{amsmath, amssymb, bm}
\usepackage{booktabs}
\usepackage[margin=1in]{geometry}

\newcommand{\R}{\mathbb{R}}
\newcommand{\C}{\mathbb{C}}
\newcommand{\tr}{\operatorname{tr}}
\newcommand{\diag}{\operatorname{diag}}
\DeclareMathOperator*{\argmin}{arg\,min}

\title{Solvers for Hierarchical PG--NB--GP:\\
  From Joint CG to Direct Schur}
\author{Scratchpad}
\date{}

\begin{document}
\maketitle

%% ================================================================
\section{Setup}

Hierarchical additive GP:
\[
  f_i = g(t_i) + h_{\ell_i}(t_i), \qquad
  g \sim \mathcal{GP}(0, k_g), \quad
  h_\ell \sim \mathcal{GP}(0, k_h), \quad \ell = 1,\dots,L.
\]
Negative binomial likelihood with P\'olya-Gamma augmentation gives the
conditional:
\[
  p(f \mid \omega, y) \propto \exp\!\bigl(-\tfrac12 f^\top \Delta\, f
  + \kappa^\top f\bigr)
  \cdot \exp\!\bigl(-\tfrac12 f^\top K^{-1} f\bigr),
\]
where $\Delta = \diag(\omega_i)$, $\kappa_i = \tfrac12(y_i - r)$,
and $K$ is the joint prior covariance of $f = g + h$.

%% ================================================================
\section{Feature-space representation}

Using the spectral approximation with shared grid $\{\xi_k\}_{k=1}^m$:
\[
  g(t) = \sum_k w_{g,k}\, a_{g,k}\, e^{2\pi i \xi_k \cdot t},
  \qquad
  h_\ell(t) = \sum_k w_{h,k}\, a_{h,k}^{(\ell)}\, e^{2\pi i \xi_k \cdot t},
\]
where $w_{g,k} = S_g(\xi_k)^{1/2}$, $w_{h,k} = S_h(\xi_k)^{1/2}$
are spectral amplitude weights.

Define:
\begin{itemize}
\item $D_g = \diag(w_{g,k}^2)$, \; $D_h = \diag(w_{h,k}^2)$,
  \; and their square roots.
\item $F \in \C^{n \times m}$: the NUFFT matrix, $F_{ik} = e^{2\pi i \xi_k t_i}$.
\item $F_\ell \in \C^{n_\ell \times m}$: restriction to site~$\ell$.
\item $T_\ell = F_\ell^* \Delta_\ell F_\ell$ (Toeplitz via NUFFT),
  \; $T_{\mathrm{all}} = \sum_\ell T_\ell$.
\end{itemize}

%% ================================================================
\section{The joint feature-space system}

The posterior precision in symmetrized feature space is the
$(1+L)m \times (1+L)m$ block matrix:
\[
  M =
  \begin{pmatrix}
    A     & B_1   & \cdots & B_L \\
    B_1^* & D_1^{(h)} &     &     \\
    \vdots &      & \ddots &     \\
    B_L^* &       &        & D_L^{(h)}
  \end{pmatrix}
\]
where:
\begin{align}
  A &= I_m + D_g^{1/2}\, T_{\mathrm{all}}\, D_g^{1/2}
     & &\text{($m \times m$, global)} \label{eq:A}\\
  D_\ell^{(h)} &= I_m + D_h^{1/2}\, T_\ell\, D_h^{1/2}
     & &\text{($m \times m$, per-site)} \label{eq:D_ell}\\
  B_\ell &= D_g^{1/2}\, T_\ell\, D_h^{1/2}
     & &\text{($m \times m$, cross-term)} \label{eq:B_ell}
\end{align}

The Schur complement for $g$:
\begin{equation}\label{eq:schur}
  S = A - \sum_{\ell=1}^L B_\ell\, (D_\ell^{(h)})^{-1}\, B_\ell^*
  \qquad (m \times m).
\end{equation}

%% ================================================================
\section{Overview of approaches}

We describe four solver approaches, ordered by development.
All share the same EM outer loop: an E-step that updates
$\tilde a_g, \tilde a_h^{(\ell)}, \sigma_{\mathrm{diag}}$,
and an M-step that updates kernel hyperparameters via
$\partial\mathcal{L}/\partial\theta$.

The key distinctions are:
\begin{enumerate}
\item How the E-step posterior mean is obtained (CG on $(1+L)m$
  system, block CD, or Schur elimination).
\item How $\sigma_{\mathrm{diag}}$ is estimated.
\item How the M-step trace term
  $\tr(M^{-1}\, \partial M/\partial\theta)$ is computed
  (observation-space Hutchinson, feature-space Hutchinson, or exact).
\end{enumerate}

%% ================================================================
\section{Approach A: Joint CG solver (baseline)}

\paragraph{E-step.}
Solve the full $(1{+}L)m$-dimensional system $M\,\tilde a = \text{rhs}$
with conjugate gradients. Each CG iteration applies the block
matvec via Toeplitz embeddings: $O((1{+}L)\,m\log m)$ per iteration.

\paragraph{sigma\_diag.}
Observation-space Hutchinson: for each probe $z \in \R^n$,
solve $(K^{-1} + \Delta)^{-1} z$ via CG on the full system,
transforming between observation and feature space at each CG step.
Requires $2(1{+}L)$ NUFFTs per CG iteration.

\paragraph{M-step.}
Observation-space Hutchinson probes: for each probe,
apply $\partial K/\partial\theta$ (Toeplitz matvec in feature space,
then NUFFT back), and compute $v^\top (\partial K/\partial\theta) v$
where $v = M^{-1} z$ from an additional CG solve.
Requires $\sim 2(1{+}L)$ NUFFTs per probe per CG iteration.

\paragraph{Cost per outer iteration.}
Dominated by NUFFTs inside CG loops. Typical CG counts:
$\sim$25,000--36,000 total CG iterations, each touching
$O(n)$ NUFFTs.

\paragraph{Benchmark.} 1331s total, 88.7s/iter (mosquito dataset,
$n{=}24{,}513$, $L{=}94$, $m{\approx}25$).

%% ================================================================
\section{Approach B: Block coordinate descent + Schur M-step}

\paragraph{Key idea.} Exploit the arrow structure of $M$:
the $h$-blocks are decoupled given $g$. Alternate between
updating $g$ (fixing all $h_\ell$) and updating each $h_\ell$
(fixing $g$), each via CG on an $m{\times}m$ system.

\paragraph{E-step.}
Each CD sweep:
\begin{itemize}
\item \textbf{g-step}: solve $A\,\tilde a_g = D_g^{1/2} F^*
  (\kappa - \Delta\,h_{\text{combined}})$ via CG.
  Cost: 1 NUFFT (type~1) for RHS, CG iterations using Toeplitz
  matvecs, 1 NUFFT (type~2) to recover $g$ in observation space.
\item \textbf{h-step}: for each $\ell$, solve
  $D_\ell^{(h)}\,\tilde a_h^{(\ell)} = D_h^{1/2} F_\ell^*
  (\kappa_\ell - \Delta_\ell g_\ell)$ via CG.
  Cost: $L$ NUFFTs (type~1) + $L$ NUFFTs (type~2).
\end{itemize}
Per sweep: $2(1{+}L)$ NUFFTs. With 5 sweeps: $10(1{+}L)$ NUFFTs.

\paragraph{sigma\_diag.}
Independent (uncoupled) Hutchinson: solve $g$ and each $h_\ell$
component separately via CG, batching probes.
Cost: $2(1{+}L)$ NUFFTs (amortized over probes).

\paragraph{M-step.}
Feature-space Schur: form $S$ from the Toeplitz blocks
(materializing $T_\ell$ via $m$ basis-vector matvecs),
then compute $\tr(S^{-1}\,\partial A/\partial\theta_g)$ and
analogous $\theta_h$ terms using Hutchinson probes in
$(1{+}L)m$-dimensional feature space. No observation-space
NUFFTs needed.

\paragraph{Cost per outer iteration.}
E-step: $3 \times [10(1{+}L) + 2(1{+}L)] = 36(1{+}L) \approx 3420$
NUFFTs. M-step: $\sim$4000 feature-space CG iterations (Toeplitz only).
No M-step NUFFTs.

\paragraph{Benchmark.} 649s total, 43.3s/iter.
Speedup: \textbf{2.0$\times$} over joint.

\paragraph{Limitation.} CD is an approximation: the $g$/$h$ blocks
are not jointly optimal at any finite number of sweeps.
Hyperparameters (especially $\sigma_g^2$) can drift relative to
the joint solve, though predictions remain accurate.

%% ================================================================
\section{Approach C: Schur CG (fully iterative)}

\paragraph{Key idea.} Eliminate $h$ analytically via the Schur
complement~\eqref{eq:schur}, reducing the E-step to an $m{\times}m$
CG solve for $\tilde a_g$. Then recover each $\tilde a_h^{(\ell)}$
by backsubstitution. This is \emph{exact} (no CD approximation)
and stays entirely in feature space after the initial NUFFT transform.

\paragraph{E-step.}
\begin{enumerate}
\item Transform $\kappa$ to feature space:
  $\hat\kappa_g = D_g^{1/2} F^* \kappa$, \;
  $\hat\kappa_{h,\ell} = D_h^{1/2} F_\ell^* \kappa_\ell$.
  Cost: $1{+}L$ NUFFTs (type~1).
\item Compute effective RHS:
  $b_g^{\text{eff}} = \hat\kappa_g
    - \sum_\ell B_\ell\,(D_\ell^{(h)})^{-1}\,\hat\kappa_{h,\ell}$.
  Each $(D_\ell^{(h)})^{-1}$ application uses CG with Toeplitz
  matvecs: $O(m\log m)$ per CG step.
\item Solve $S\,\tilde a_g = b_g^{\text{eff}}$ via CG.
  Each CG matvec applies $S = A - \sum_\ell B_\ell (D_\ell^{(h)})^{-1} B_\ell^*$,
  requiring $L$ inner CG solves of size $m$ per outer CG step.
  All Toeplitz, no NUFFTs.
\item Recover $\tilde a_h^{(\ell)}
  = (D_\ell^{(h)})^{-1}(\hat\kappa_{h,\ell} - B_\ell^* \tilde a_g)$.
\item Transform back: $g = F\,D_g^{1/2}\,\tilde a_g$ etc.
  Cost: $1{+}L$ NUFFTs (type~2).
\end{enumerate}
Total: $2(1{+}L)$ NUFFTs per E-step inner iteration.

\paragraph{sigma\_diag.} Same as Approach~B: independent batched
Hutchinson with $2(1{+}L)$ NUFFTs.

\paragraph{M-step.} Feature-space Hutchinson probes with the
Schur matvec (same as Approach~B). No NUFFTs.

\paragraph{Cost per outer iteration.}
E-step: $3 \times [2(1{+}L) + 2(1{+}L)] = 12(1{+}L) \approx 1140$
NUFFTs (3$\times$ fewer than Block~CD). But the Schur CG matvec
is expensive: each application requires $L$ inner CG solves,
so the Toeplitz matvec count is high.

\paragraph{Benchmark.} 619s total, 41.3s/iter.
Speedup: \textbf{2.2$\times$} over joint.

\paragraph{Tradeoff.} Fewer NUFFTs than Block~CD, but the nested
CG (inner solves inside outer CG) has high Toeplitz overhead.
Generalizes naturally to $d \geq 2$ since no matrix factorization
is needed.

%% ================================================================
\section{Approach D: Direct Schur (dense factorization)}

\paragraph{Key idea.} Since $m \approx 25$ (for $d{=}1$ with
typical lengthscales), explicitly form $T_\ell$, $A$, $D_\ell^{(h)}$,
$B_\ell$, and $S$ as dense $m{\times}m$ matrices, then use Cholesky
factorization for all solves. This eliminates CG entirely from the
E-step and M-step.

\paragraph{Block construction.}
Form each $T_\ell \in \C^{m \times m}$ by applying the Toeplitz
embedding to $m$ standard basis vectors: cost $O(m^2 \log m)$ per site.
Then:
\begin{align}
  A &= I + D_g^{1/2}\, T_{\mathrm{all}}\, D_g^{1/2}, &
  &\text{Cholesky: } O(m^3) \\
  D_\ell^{(h)} &= I + D_h^{1/2}\, T_\ell\, D_h^{1/2}, &
  &\text{Cholesky: } O(Lm^3) \\
  B_\ell &= D_g^{1/2}\, T_\ell\, D_h^{1/2}, &
  &O(Lm^2) \\
  S &= A - \textstyle\sum_\ell B_\ell\, (D_\ell^{(h)})^{-1}\, B_\ell^*, &
  &\text{Cholesky: } O(Lm^3)
\end{align}
Total: $O(Lm^3)$. With $L{=}94$, $m{=}25$: $\sim 1.5 \times 10^6$
ops. \textbf{Negligible} ($< 1$\,ms).

\paragraph{E-step.}
\begin{enumerate}
\item Transform $\kappa$ to feature space: $1{+}L$ NUFFTs.
\item Compute effective RHS using pre-factorized
  $(D_\ell^{(h)})^{-1}$: direct backsolves, $O(Lm^2)$.
\item Solve $S\,\tilde a_g = b_g^{\text{eff}}$: Cholesky backsolve,
  $O(m^2)$.
\item Recover $\tilde a_h^{(\ell)}$: direct backsolves, $O(Lm^2)$.
\item Transform back: $1{+}L$ NUFFTs.
\end{enumerate}
Total: $2(1{+}L)$ NUFFTs + $O(Lm^2)$ algebra per inner iteration.

\paragraph{sigma\_diag.}
\textbf{Exact} Schur diagonal: compute $[M^{-1}]_{ii}$ for each
observation using the Schur block inversion formulas and
pre-factorized matrices. Can also use Hutchinson with direct
solves (no CG). Cost: $2(1{+}L)$ NUFFTs.

\paragraph{M-step.}
\textbf{Exact trace} --- no Hutchinson probes needed.
Since $S^{-1}$ and $\partial A/\partial\theta$ are available as
dense $m{\times}m$ matrices:
\begin{align}
  \tr\bigl(M^{-1}\, \partial M / \partial\theta_g\bigr)
  &= \tr\bigl(S^{-1}\, \partial A / \partial\theta_g\bigr), &
  &O(m^2) \\
  \tr\bigl(M^{-1}\, \partial M / \partial\theta_h\bigr)
  &= \sum_\ell \bigl[\tr(\cdots) + 2\Re\,\tr(\cdots)\bigr], &
  &O(Lm^2)
\end{align}
\textbf{Zero NUFFTs, zero CG, zero stochastic estimation.}

\paragraph{Cost per outer iteration.}
E-step: $3 \times 2(1{+}L) = 570$ NUFFTs + $O(Lm^2)$ dense algebra.
M-step: $O(Lm^2)$ exact. \textbf{No CG anywhere.}

\paragraph{Benchmark.} 159s total, 10.6s/iter.
Speedup: \textbf{8.4$\times$} over joint.

\paragraph{Limitation.} Requires $m$ to be small enough for dense
$m{\times}m$ factorization ($m \lesssim 200$). This holds for $d{=}1$
with typical lengthscales, but for $d \geq 2$ the spectral grid
grows as $m^d$, making dense factorization infeasible. In that regime,
Approach~C (Schur CG) is preferred.

%% ================================================================
\section{Comparison}

\begin{table}[h]
\centering
\begin{tabular}{@{}lcccc@{}}
\toprule
& \textbf{A: Joint} & \textbf{B: Block CD} & \textbf{C: Schur CG} & \textbf{D: Direct} \\
\midrule
E-step solve & CG on $(1{+}L)m$ & CD sweeps, CG on $m$ & Schur CG on $m$ & Cholesky on $m$ \\
E-step NUFFTs/inner & many (in CG) & $10(1{+}L)$ & $2(1{+}L)$ & $2(1{+}L)$ \\
sigma\_diag & obs-space CG & indep.\ CG & indep.\ CG & exact / direct \\
M-step & obs-space Hutch. & feat-space Hutch. & feat-space Hutch. & exact trace \\
M-step NUFFTs & many & 0 & 0 & 0 \\
CG required? & yes & yes & yes (nested) & \textbf{no} \\
Exact joint post.? & yes & no (CD approx.) & yes & yes \\
Scales to $d{\geq}2$? & yes & yes & yes & no ($m^d$) \\
\midrule
Time/iter (s) & 88.7 & 43.3 & 41.3 & \textbf{10.6} \\
Total time (s) & 1331 & 649 & 619 & \textbf{159} \\
Speedup vs A & 1.0$\times$ & 2.0$\times$ & 2.2$\times$ & \textbf{8.4$\times$} \\
Final MAE & --- & 27.4 & 27.3 & 27.4 \\
\bottomrule
\end{tabular}
\caption{Comparison on mosquito dataset ($n{=}24{,}513$, $L{=}94$, $m{\approx}25$,
  15 outer iterations, 3 inner E-steps).}
\end{table}

\paragraph{Converged hyperparameters.}
\begin{table}[h]
\centering
\begin{tabular}{@{}lccccc@{}}
\toprule
& $\ell_g$ & $\sigma_g^2$ & $\ell_h$ & $\sigma_h^2$ & $r$ \\
\midrule
B: Block CD & 0.085 & 4.04 & 0.105 & 2.10 & 2.84 \\
C: Schur CG & 0.196 & 1.57 & 0.072 & 1.17 & 2.84 \\
D: Direct & 0.313 & 0.95 & 0.083 & 1.85 & 2.84 \\
\bottomrule
\end{tabular}
\caption{Learned hyperparameters after 15 iterations.
  Despite different $\ell_g / \sigma_g^2$ tradeoffs, all approaches
  reach similar MAE (${\approx}27.4$), indicating that the posterior
  predictive is robust to this identifiability.}
\end{table}

%% ================================================================
\section{NUFFT budget analysis}

NUFFTs are $O(n)$ and dominate wall-clock time. Toeplitz matvecs
are $O(m\log m)$ and essentially free for $m{=}25$.

\begin{table}[h]
\centering
\begin{tabular}{@{}lrrr@{}}
\toprule
& \textbf{B: Block CD} & \textbf{C: Schur CG} & \textbf{D: Direct} \\
\midrule
E-step (per inner) & $10(1{+}L)=950$ & $2(1{+}L)=190$ & $2(1{+}L)=190$ \\
sigma\_diag (per inner) & $2(1{+}L)=190$ & $2(1{+}L)=190$ & $2(1{+}L)=190$ \\
Inner iters $\times 3$ & 3420 & 1140 & 1140 \\
M-step & 0 & 0 & 0 \\
\midrule
\textbf{Total per outer} & \textbf{3420} & \textbf{1140} & \textbf{1140} \\
\bottomrule
\end{tabular}
\caption{NUFFT calls per outer iteration. Approaches C and D have identical
NUFFT budgets; D is faster because it replaces CG with direct solves.}
\end{table}

%% ================================================================
\section{Recommendations}

\begin{itemize}
\item \textbf{$d{=}1$, moderate $m$}: Use Approach~D (Direct Schur).
  8.4$\times$ faster than joint, no CG, exact traces.
  Works whenever $m \lesssim 200$.

\item \textbf{$d \geq 2$ or very small lengthscale}:
  Use Approach~C (Schur CG).
  Same NUFFT budget as Direct but avoids dense $m^d {\times} m^d$
  factorization. The nested CG adds Toeplitz overhead but
  scales gracefully.

\item \textbf{Approach B} (Block CD) is superseded: it uses
  5$\times$ more NUFFTs than Schur-based approaches and
  introduces a CD approximation error in the joint posterior.
\end{itemize}

\end{document}
